\section{Activation layer}\label{sec:act_appendix}


\begin{definition}[Auxiliary activation]\label{def:act}
    For a continuous function $\act: \RR \to \RR$,  an activation layer takes $\vx\in\RR^{\aux{D}}$ as input and outputs $\vy = \act(\vx)$ with $ y_i = \act (x_i) $ for all $i \le \aux{D}$.
\end{definition}

  In the discussions below, we consider an activation layer in the auxiliary model with activation function $\act$ that takes in input sequence $\vx_1, \cdots, \vx_{\aux{T}}$ and outputs $\vy_1, \cdots, \vy_{\aux{T}}$, with $\vy_t =\act(\vx_t)$ for each $t \le {\aux{T}}$. Since this involves a token-wise operation, we will present our constructed modules with a general token position $t$. Since no parameters of the auxiliary model are involved in this operation, the prefix tokens $\{ \vv_j \}$ contain $0$ in the following modules.

\paragraph{\simulator{} Activation Forward module} The embedding $\ve_t$ contains $\vx_t$ in its first $\aux{D}$ indices. We simply pass the embeddings into activation $\act$, which returns $\act(\vx_t)$ in its first $\aux{D}$ indices.



\paragraph{Auxiliary's backpropagation through activation} With the definition in \cref{def:act}, the auxiliary's backpropagation takes in the loss gradient w.r.t. output ($\losspartial{\vy}$) and computes the loss gradient  w.r.t. input ($\losspartial{\vx}$). We further assume that the derivative of $\act$ is well-defined everywhere. This assumption includes non-differentiable activation functions with well-defined derivatives like $ReLU$.

\begin{definition}[Auxiliary activation backpropagation]\label{def:act_backprop}
   For a continuous function $\act: \RR \to \RR$, with a well-defined derivative $\act'(x) = \partial \act(x) / \partial x$ for each $x \in \RR$,  the backpropagation takes $\losspartial{ \vy }, \vx \in\RR^{\aux{D}}$ as input and outputs
    \begin{align*}
        \losspartial{ \vx } = \act'(\vx) \odot \losspartial{\vy},
    \end{align*}
    where $\act'(\vx) \in \RR^{\aux{D}}$ with $\act'(\vx)_i = \act'(x_i)$ at each $i \le \aux{D}.$  
\end{definition}

\paragraph{Complexity of true backpropagation} The above operation is computation heavy since 
 it involves $\act'(\vx) \odot \losspartial{\vy}$. As mentioned for the layer normalization module, the element-wise multiplication between $\act'(\vx)$ and $\losspartial{\vy}$ will require an MLP module following \cref{thm:gelu_multiplication}. Furthermore, it involves changing the activation function in \simulator{} in specific modules to $\act'$. To circumvent this, we instead turn to a first-order Taylor approximation.


\begin{definition}[Approximate Activation backpropagation]\label{def:approx_act_backprop}
   For a continuous function $\act: \RR \to \RR$ and a hyperparameter $\epsilon$, the layer takes $\losspartial{ \vy }, \vx \in\RR^{\aux{D}}$ as input and outputs
    \begin{align*}
        \Hat{ \losspartial{ \vx } } = \frac{1}{\epsilon} \left( \act(\vx + \epsilon \losspartial{\vy}) - \act(\vx) \right).
    \end{align*}
\end{definition}



The following theorems show that under mild assumptions on the activation function and the input, gradient pair, the first-order gradient is a good approximation to the true gradient.


\begin{restatable}{theorem}{firstorderactbckprop}\label{thm:first_order_act_backprop}    
%\end{restatable}
%\begin
For any $\epsilon > 0$, $B_y, B_{act} > 0$, consider a second-order differentiable activation function $\act: \RR \to \RR$, with $\partial^2 \act (x) / \partial (x^2) $ bounded by $B_{act}$ for each $x \in \RR$. Then, for any  input $\vx \in \RR^{\aux{D}}$ and gradient $\losspartial{\vy} \in \RR^{\aux{D}}$ with $\norm{\losspartial{\vy}}_{2} \le B_y$, the following holds true:
\begin{align*}
    \norm{ \losspartial{ \vx } - \Hat{ \losspartial{ \vx } } }_{2} \le \mathcal{O}( B_{act} B_y^2 \epsilon  ),
\end{align*}
where $\losspartial{ \vx }, \Hat{ \losspartial{ \vx } }$ have been defined using $\vx, \losspartial{\vy}$, and $\epsilon$ in \cref{def:act_backprop,def:approx_act_backprop}.
%\end{theorem}
\end{restatable}


For $\mathrm{ReLU}$ activation, which is not second-order differentiable at $0$, we instead bound the difference between $\losspartial{ \vx }, \Hat{ \losspartial{ \vx } }$ by defining some form of alignment between input and  gradient pair $\vx, \losspartial{\vy}$.

\begin{definition}[$(\epsilon, \rho)$-alignment]\label{def:eps-rho_aligned}
Input and gradient $\vx, \losspartial{\vy} \in \RR^{\aux{D}}$ are said to be $(\epsilon, \rho)$-aligned, if there exist a set $C \subseteq [\aux{D}]$, with  $\abs{C} \ge (1 - \rho) \aux{D}$, such that for each $i$ in $C$, $\abs{x_i} > \epsilon \abs{ (\losspartial{\vy} )_i }.$
\end{definition}

$\epsilon$ controls the fraction of coordinates where $\abs{x_i} \leq \epsilon \abs{ (\losspartial{\vy} )_i }$. As $\epsilon \to 0$, $\rho \to 0$  as well for bounded gradients. 

%\begin{example}
\begin{restatable}{example}{examplealignment}
    For any $B_{min}, B_{max} > 0$, all inputs $\vx$ that satisfy $\min_i \abs{x_i} > B_{min}$ , and gradients $\losspartial{\vy}$ that satisfy $\max_j \abs{(\losspartial{\vy})_j} \le B_{max}$, are $( B_{min} / B_{max}, 0)$-aligned.
\end{restatable}
    
%\end{example}


%\begin{theorem}\label{thm:first_order_act_backprop_relu}
\begin{restatable}{theorem}{firstorderactbckproprelu}\label{thm:first_order_act_backprop_relu}
For any $\epsilon, \rho > 0$ and $B_y > 0$, for any input $\vx \in \RR^{\aux{D}}$ and gradient $\losspartial{\vy} \in \RR^{\aux{D}}$, with $\norm{\losspartial{\vy}}_{\infty} \le B_y$,  that are $(\epsilon, \rho)$-aligned by \cref{def:eps-rho_aligned}, 
\begin{align*}
    \norm{ \losspartial{ \vx } - \Hat{ \losspartial{ \vx } } }_{2} \le \mathcal{O}( B_y \sqrt{\rho \aux{D}} ).
\end{align*}
where $\losspartial{ \vx }, \Hat{ \losspartial{ \vx } }$ have been defined using $\vx, \losspartial{\vy}$, $\epsilon$ and $\act = \mathrm{ReLU}$ in \cref{def:act_backprop,def:approx_act_backprop}.
%\end{theorem}
\end{restatable}


%For example, for all inputs $\vx$ that additionally satisfy $\min_i \abs{x_i} \ge B_{x}$ for some $B_x > 0$, we can set $\epsilon \leq \frac{B_y}{B_x}$

\paragraph{\simulator{} Activation backpropagation module}
The input embeddings contain $\losspartial{\vy_t}$ in the first $\aux{D}$ embeddings. 
With the requirement of the activation layer input for gradient, we copy $\vx_t$ from the Forward module at each position $t$. We set $\epsilon$ as a hyper-parameter and return $\Hat{ \losspartial{ \vx_t } }$ as the output of this module.


$\Hat{ \losspartial{ \vx_t } }$ will be computed using a single-layer MLP with activation $\act$ as follows. The first linear layer of the MLP will be used to compute $\vx_t + \epsilon \losspartial{\vy_t}$ and $\vx_t$. After the activation $\act$, the embedding $\ve_t$ contains $\act(\vx_t + \epsilon \losspartial{\vy_t})$ and $\act(\vx_t)$. The final linear layer of the MLP will be used to compute $\frac{1}{\epsilon} \left( \act(\vx_t + \epsilon \losspartial{\vy_t}) - \act(\vx_t) \right)$.


\subsection{Proofs of theorems}
We restate the theorems, before presenting their proofs for easy referencing.


\firstorderactbckprop*
\begin{proof}
    The proof follows along the lines of \cref{thm:first_order_LN_backprop}. Recall that given an input $\vx$, the activation layer outputs $\vy = \act(\vx)$, where the function $\act$ is applied coordinate-wise on $\vx$. Given input $\vx$ and the output gradient $\losspartial{\vy}$, the gradient w.r.t. the input is given by $\losspartial{\vx} = \act'(\vx) \odot \losspartial{\vy}$, where the $\act'$ function is also applied coordinate wise to $\vx$. We defined $\Hat{\losspartial{\vx}}$ as an $\epsilon$-approximate gradient, given by $\frac{1}{\epsilon} ( \act(\vx + \epsilon \losspartial{\vy}) - \act(\vx ) ).$ Since both $\act$ and $\act'$ are applied coordinate-wise, we can look at the coordinate-wise difference between $\losspartial{\vx}$ and $\Hat{\losspartial{\vx}}$.

    Consider an arbitrary coordinate $i \le \aux{D}$. Under the assumption that $\act$ is second-order differentiable, we have
    \begin{align*}
       (\Hat{\losspartial{\vx}})_i &= \frac{1}{\epsilon} \left( \act( x_i + \epsilon (\losspartial{\vy})_i ) - \act( x_i) \right) \\&
        = \act'(x_i)  (\losspartial{\vy})_i + \frac{1}{\epsilon} \int_{\theta=0}^{\epsilon} \frac{\partial^2 \act( x_{\theta} ) }{ \partial x^2_{\theta} } (\losspartial{\vy})^2_i  \theta d\theta \\&
        = \act'(x_i)  (\losspartial{\vy})_i + \mathcal{O}( \epsilon B_{act} (\losspartial{\vy})_i^2 ),
    \end{align*}
    where $x_{\theta}$ represents $ x_i + \theta (\losspartial{\vy})_i$ in the second step. In the final step, we utilize the upper bound assumption on $\frac{\partial^2 \act(x)}{ \partial x^2 }.$

    Thus, $(\losspartial{\vx})_i - (\Hat{\losspartial{\vx}})_i = \mathcal{O}( \epsilon B_{act} (\losspartial{\vy})_i^2 )$, and so
    \begin{align*}
        \norm{ \losspartial{\vx} - \Hat{\losspartial{\vx}} }_2 = \mathcal{O} ( \epsilon B_{act} \sum_{i=1}^{\aux{D}}(\losspartial{\vy})_i^2 ) = \mathcal{O} ( \epsilon B_{act} \norm{ \losspartial{\vy} }_2^2 ) \le \mathcal{O} ( \epsilon B_{act}  B_y^2 ).
    \end{align*}
\end{proof}



\examplealignment*
\begin{proof}
    Recall the definition of $(\epsilon, \rho)$-alignment from \cref{def:eps-rho_aligned}. Input and gradient $\vx, \losspartial{\vy} \in \RR^{\aux{D}}$ are said to be $(\epsilon, \rho)$-aligned, if there exist a set $C \subseteq [\aux{D}]$, with  $\abs{C} \ge (1 - \rho) \aux{D}$, such that for each $i$ in $C$, $\abs{x_i} > \epsilon \abs{ (\losspartial{\vy} )_i }.$ 

    Consider an arbitrary coordinate $i \le \aux{D}$. We have $\abs{x_i} > \epsilon  \abs{(\losspartial{\vy})_i}$ for any $\epsilon < \abs{x_i} / \abs{(\losspartial{\vy})_i}$. Under the assumption that $\abs{x_i} > B_{min}$, and $\abs{(\losspartial{\vy})_i} \le B_{max}$, a bound of $B_{min} / B_{max}$ suffices. 
\end{proof}

\firstorderactbckproprelu*
\begin{proof}
    Recall that given an input $\vx$, the activation layer outputs $\vy = \act(\vx)$, where the function $\act$ is applied coordinate-wise on $\vx$. Given input $\vx$ and the output gradient $\losspartial{\vy}$, the gradient w.r.t. the input is given by $\losspartial{\vx} = \act'(\vx) \odot \losspartial{\vy}$, where the $\act'$ function is also applied coordinate wise to $\vx$. We defined $\Hat{\losspartial{\vx}}$ as an $\epsilon$-approximate gradient, given by $\frac{1}{\epsilon} ( \act(\vx + \epsilon \losspartial{\vy}) - \act(\vx ) ).$ Since both $\act$ and $\act'$ are applied coordinate-wise, we can look at the coordinate-wise difference between $\losspartial{\vx}$ and $\Hat{\losspartial{\vx}}$. For $\mathrm{ReLU}$ activation, $\act'(x) = \mathrm{sign}(x)$ for all $x \in \RR \setminus \{0\}$, with $\act'(0) = 1$ to avoid ambiguity. 
    
    Going by the definition of $(\epsilon, \rho)$-alignment of the input and gradient from \cref{def:eps-rho_aligned}, we have a set $C$ with $\abs{C} \ge (1-\rho) \aux{D}$ such that for each $i \in \aux{D}$, $\abs{x_i} > \epsilon  \abs{(\losspartial{\vy})_i}$. For all coordinates $i \in C$, we can then observe that $\mathrm{sign} (x_i + \epsilon (\losspartial{\vy})_i) = \mathrm{sign} (x_i)$, implying 
    $$
    \act( x_i + \epsilon (\losspartial{\vy})_i ) - \act( x_i ) = \epsilon (\losspartial{\vy})_i \act'(x_i) = \epsilon (\losspartial{\vx})_i
    $$

    
    For coordinates $i \notin C$, we have three possible cases:
    \begin{itemize}
        \item $\mathrm{sign} (x_i) = \mathrm{sign} (x_i + \epsilon (\losspartial{\vy})_i ) $: In this case, we can again show $\act( x_i + \epsilon (\losspartial{\vy})_i ) - \act( x_i ) = \epsilon (\losspartial{\vy})_i \act'(x_i) = \epsilon (\losspartial{\vx})_i.$

        \item $\mathrm{sign} (x_i) = 0$, $\mathrm{sign} (x_i + \epsilon (\losspartial{\vy})_i ) = 1$: In this case, we have $\act'(x_i) = 0$, and so $(\losspartial{\vx})_i = 0$. Additionally, $\mathrm{sign} ( (\losspartial{\vy})_i ) = 1$, and so
        \begin{align*}
        \abs{ \act( x_i + \epsilon (\losspartial{\vy})_i ) - \act( x_i ) - \epsilon (\losspartial{\vx})_i } &= \abs{ x_i + \epsilon (\losspartial{\vy})_i} 
        \leq \epsilon \abs{  (\losspartial{\vy})_i },
        \end{align*}
        where in the final step, we use the fact that $x_i < 0$ and $\abs{x_i} < \epsilon \abs{ (\losspartial{\vy})_i }.$ 
        
        \item  $\mathrm{sign} (x_i) = 1$, $\mathrm{sign} (x_i + \epsilon (\losspartial{\vy})_i ) = 0$: In this case, we have $\act'(x_i) = 1$, and so $(\losspartial{\vx})_i = (\losspartial{\vy})_i$.  Additionally, $\mathrm{sign} ( (\losspartial{\vy})_i ) = 0$, and so
        \begin{align*}
        \abs{ \act( x_i + \epsilon (\losspartial{\vy})_i ) - \act( x_i ) - \epsilon (\losspartial{\vx})_i } &= \abs{ - x_i  - \epsilon (\losspartial{\vy})_i } 
        \leq \abs{ \epsilon (\losspartial{\vy})_i },
        \end{align*}
        where in the final step, we use the fact that $x_i \ge 0$ and $\abs{x_i} < \epsilon \abs{ (\losspartial{\vy})_i }.$ 
        
    \end{itemize}

    Thus, from the above discussion, we have
    \begin{align*}
        \norm{ \losspartial{\vx} - \Hat{\losspartial{\vx}} }_2 &= \frac{1}{\epsilon} \left( \sum_{i=1}^{\aux{D}} (\act( x_i + \epsilon (\losspartial{\vy})_i ) - \act( x_i ) - \epsilon (\losspartial{\vx})_i)^2 \right)^{1/2} \\&
        = \frac{1}{\epsilon} \left(  \sum_{i \notin C} (\act( x_i + \epsilon (\losspartial{\vy})_i ) - \act( x_i ) - \epsilon (\losspartial{\vx})_i)^2 \right)^{1/2} \\&
        \leq \left(  \sum_{i \notin C}  (\losspartial{\vy})_i^2 \right)^{1/2} \leq \sqrt{\rho \aux{D}} \sqrt{\max_{i \notin C}  (\losspartial{\vy})^2_i } \leq \sqrt{\rho \aux{D}} B_y.
    \end{align*}
    The final step includes a simple Cauchy Schwartz inequality and the desired bound comes from the assumed bound on $\norm{ \losspartial{\vy} }_2 $.
\end{proof}